\par\medskip
\pdfbookmark[4]{5.}{tabula-nova-005}
\begin{center}5.\end{center}
\noindent\hspace*{\parindent}Statuendo argumentum partis tertiae tabulae nostrae, puta quantitatem
\[
\BareRoot[3]{\frac{B}{t}}=206265''\mathpunct{.}\beta=\frac{648000''}{\pi}\beta,
\]
habetur
\[
\theta+\tfrac13\theta^3=\frac{8}{3\beta^3}\mathpunct{\raisebox{.45ex}{,}}
\]
\SourceLeaf{361}quam formulam, statuendo \(\dfrac{1}{\theta^3}=\eta\), ita exhibere possumus
\[
\eta=\frac{\beta^3}{8}+\frac{3\beta^3}{8}\eta^{\frac23}.
\]
Applicando huic aequationi theorema \textsc{Lagrang}ianum, invenimus generaliter
\[
\begin{aligned}
\eta^{\frac n3}=\frac{1}{\theta^n}={}&\left(\frac{\beta}{2}\right)^n+n\left(\frac{\beta}{2}\right)^{n+2}
+\frac{n\mathpunct{.}n+1}{1\mathpunct{.}2}\left(\frac{\beta}{2}\right)^{n+4}\\
&+\frac{n\mathpunct{.}n\mathpunct{.}n+3}{1\mathpunct{.}2\mathpunct{.}3}\left(\frac{\beta}{2}\right)^{n+6}\\
&+\frac{n\mathpunct{.}n-1\mathpunct{.}n+2\mathpunct{.}n+5}{1\mathpunct{.}2\mathpunct{.}3\mathpunct{.}4}\left(\frac{\beta}{2}\right)^{n+8}\\
&+\frac{n\mathpunct{.}n-2\mathpunct{.}n+1\mathpunct{.}n+4\mathpunct{.}n+7}{1\mathpunct{.}2\mathpunct{.}3\mathpunct{.}4\mathpunct{.}5}\left(\frac{\beta}{2}\right)^{n+10}+\text{etc.,}
\end{aligned}
\]
cuius seriei lex facile elucet. Hinc fit
\[
\begin{aligned}
2\operatorname{arc.}\operatorname{tang}\frac1\theta
={}&\frac2\theta-\frac2{3\theta^3}+\frac2{5\theta^5}-\frac2{7\theta^7}+\text{etc.}\\
={}&\beta-\tfrac13\cdot\tfrac14\left(1-\tfrac31\right)\beta^3
+\tfrac15\cdot\tfrac1{16}\left(1-\tfrac51+\tfrac{5\mathpunct{.}2}{1\mathpunct{.}2}\right)\beta^5\\
&-\tfrac17\cdot\tfrac1{64}\left(1-\tfrac71+\tfrac{7\mathpunct{.}4}{1\mathpunct{.}2}-\tfrac{7\mathpunct{.}4\mathpunct{.}1}{1\mathpunct{.}2\mathpunct{.}3}\right)\beta^7\\
&+\tfrac19\cdot\tfrac1{256}\Bigl(1-\tfrac91+\tfrac{9\mathpunct{.}6}{1\mathpunct{.}2}-\tfrac{9\mathpunct{.}6\mathpunct{.}3}{1\mathpunct{.}2\mathpunct{.}3}\\
&\hspace{5em}+\tfrac{9\mathpunct{.}6\mathpunct{.}3\mathpunct{.}0}{1\mathpunct{.}2\mathpunct{.}3\mathpunct{.}4}\Bigr)\beta^9-\text{etc.}
\end{aligned}
\]
\par
Haec series exhibet complementum anomaliae ad \(180^0\) in partibus radii expressam, coëfficentiumque lex satis manifesta est, scilicet coëfficiens potestatis \(\beta^{2m+1}\) erit generaliter
\[
\begin{aligned}
&\pm\frac{1}{2m+1}\times\frac{1}{2^{2m}}\times\text{aggregatum }m+1\text{ terminorum}\\
&\text{primorum seriei, quae ex evolutione binomii }(1-3)^{\frac{2m+1}{3}}\text{ oritur.}
\end{aligned}
\]
Eruendo valores numericos, habemus in partibus radii
\[
\begin{aligned}
v={}&\pi-\beta-\tfrac16\beta^3-\tfrac1{80}\beta^5+\tfrac5{672}\beta^7+\tfrac1{288}\beta^9\\
&+\tfrac7{16896}\beta^{11}-\tfrac{113}{479232}\beta^{13}-\tfrac1{7680}\beta^{15}\text{etc.,}
\end{aligned}
\]
adeoque in minutis secundis
\[
v=180^0-\BareRoot[3]{\frac{B}{t}}-\tfrac16\frac{B}{t}\frac1{206265^2}
-\tfrac1{80}\left(\frac{B}{t}\right)^{\frac53}\frac1{206265^4}+\text{etc.}
\]
\par
Hinc perspicitur ratio propter quam quantitatem \(\BareRoot[3]{\frac{B}{t}}\) potissimum tamquam argumentum tabulae in parte tertia adoptavimus; ita scilicet argumentum proxime cum complemento anomaliae ad \(180^0\) convenit.
\par
Pro radio vectore nanciscimur, statuendo \(n=-2\),
\SourceLeaf{362}
\[
\begin{aligned}
r={}&q\Biggl(\left(\frac2\beta\right)^2-1+\left(\frac\beta2\right)^2
+\frac{-2\mathpunct{.}-2\mathpunct{.}+1}{1\mathpunct{.}2\mathpunct{.}3}\left(\frac\beta2\right)^4+\ast\\
&+\frac{-2\mathpunct{.}-4\mathpunct{.}-1\mathpunct{.}+2\mathpunct{.}+5}{1\mathpunct{.}2\mathpunct{.}3\mathpunct{.}4\mathpunct{.}5}\left(\frac\beta2\right)^8\\
&+\frac{-2\mathpunct{.}-5\mathpunct{.}-2\mathpunct{.}+1\mathpunct{.}+4\mathpunct{.}+7}{1\mathpunct{.}2\mathpunct{.}3\mathpunct{.}4\mathpunct{.}5\mathpunct{.}6}\left(\frac\beta2\right)^{10}+\ast+\text{etc.}\Biggr)
\end{aligned}
\]
\[
\begin{aligned}
={}&q\Biggl(\left(\frac2\beta\right)^2-1+\left(\frac\beta2\right)^2
+\frac{1\mathpunct{.}2\mathpunct{.}2}{1\mathpunct{.}2\mathpunct{.}3}\left(\frac\beta2\right)^4+\ast\\
&-\frac{1\mathpunct{.}2\mathpunct{.}2\mathpunct{.}4\mathpunct{.}5}{1\mathpunct{.}2\mathpunct{.}3\mathpunct{.}4\mathpunct{.}5}\left(\frac\beta2\right)^8\\
&-\frac{1\mathpunct{.}2\mathpunct{.}2\mathpunct{.}4\mathpunct{.}5\mathpunct{.}7}{1\mathpunct{.}2\mathpunct{.}3\mathpunct{.}4\mathpunct{.}5\mathpunct{.}6}\left(\frac\beta2\right)^{10}+\ast\\
&+\frac{1\mathpunct{.}2\mathpunct{.}2\mathpunct{.}4\mathpunct{.}5\mathpunct{.}7\mathpunct{.}8\mathpunct{.}11}{1\mathpunct{.}2\mathpunct{.}3\mathpunct{.}4\mathpunct{.}5\mathpunct{.}6\mathpunct{.}7\mathpunct{.}8}\left(\frac\beta2\right)^{14}\\
&+\frac{1\mathpunct{.}2\mathpunct{.}2\mathpunct{.}4\mathpunct{.}5\mathpunct{.}7\mathpunct{.}8\mathpunct{.}10\mathpunct{.}13}{1\mathpunct{.}2\mathpunct{.}3\mathpunct{.}4\mathpunct{.}5\mathpunct{.}6\mathpunct{.}7\mathpunct{.}8\mathpunct{.}9}\left(\frac\beta2\right)^{16}+\ast\\
&-\frac{1\mathpunct{.}2\mathpunct{.}2\mathpunct{.}4\mathpunct{.}5\mathpunct{.}7\mathpunct{.}8\mathpunct{.}10\mathpunct{.}11\mathpunct{.}14\mathpunct{.}17}{1\mathpunct{.}2\mathpunct{.}3\mathpunct{.}4\mathpunct{.}5\mathpunct{.}6\mathpunct{.}7\mathpunct{.}8\mathpunct{.}9\mathpunct{.}10\mathpunct{.}11}\left(\frac\beta2\right)^{20}-\text{etc.}\Biggr)
\end{aligned}
\]
sive
\[
\begin{aligned}
[r]={}&q\Biggl[\Biggl(\left(\frac2\beta\right)^2-1+\left(\frac\beta2\right)^2
+\frac23\left(\frac\beta2\right)^4-\frac23\left(\frac\beta2\right)^8\\
&-\frac79\left(\frac\beta2\right)^{10}+\frac{11}{9}\left(\frac\beta2\right)^{14}+\frac{130}{81}\left(\frac\beta2\right)^{16}
-\frac{238}{81}\left(\frac\beta2\right)^{20}-\ldots\Biggr)\Biggr].
\end{aligned}
\]
\par
Ceterum utraque series converget quamdiu \(\beta<\BareRoot[3]{4}\) sive quamdiu anomalia vera maior est [quam] \(61^0\ 35'\ 45''{,}4\).
\par\medskip
\begin{center}\rule{.18\textwidth}{.35pt}\end{center}
